% page 225 of the facsimile (image p-224 of the scan)
% block 167.6 x 237.7 mm ; left margin 34.4 mm
\pgc{12.700mm}{0.000mm}{
\l{\cel{56}217}
\l{}
\l{\sou{higrometro}.\cel{2}Aparato per qua onu mezuras la quanto de aquo-}
\l{\cel{3}vaporo qua existas en atmosfero, per korpo quan humideso}
\l{\cel{3}igas modifikesar pezale, voluminale. - DEFIRS.}
\l{}
\l{\sou{higrometrio}. Parto di fiziko, qua mezuras la quanto de aquo-}
\l{\cel{3}vaporo quan kontenas atmosfero. - DEFIRS.}
\l{}
\l{\sou{higroskopo}. Aparato per qua onu konstatas la humideso-grado di}
\l{\cel{3}atmosfero. - DEF.}
\l{}
\l{\sou{hike}. I. En la loko ube esas la persono qua pronuncas ca vorto.}
\l{\cel{3}- II. En ica loko di la dico. - L.}
\l{}
\l{\sou{hilo.}\cel{2}I. (\sou{bot.}) Marko di la ligo-punto di la grano sur la fu-}
\l{\cel{3}nikulo. - II. (\sou{anat.}) Ligo-punto ube vaskulo juntesas a vicero.}
\l{\cel{3}- EFIS.}
\l{}
\l{\sou{hiloto}. I. (\sou{epoki antiqua})\cel{2}Sklavo, che la Spartani. - II. La}
\l{\cel{3}viro, en socio, qua esas sur la grado maxim infra di la abjekt-}
\l{\cel{3}eso-skalo. - DEFIRS.}
\l{}
\l{\sou{himeno}. I. (\sou{anat.}) Membrano qua klozas parte la vagino, che la}
\l{\cel{3}virgini. - II. (\sou{bot.}) Pelikulo di la korolo qua laceresas lor}
\l{\cel{3}la desklozeso. - DEFIS.}
\l{}
\l{\sou{himeneo}. Kanto mariajala, che la Greki antiqua, e che la Romani}
\l{\cel{3}lora. - DEFIRS.}
\l{}
\l{\sou{himenio}. (\sou{bot.}) Che la fungo, la mikra membrano en qua esas la}
\l{\cel{3}spori.}
\l{}
\l{\sou{himenoptero}. (\sou{zool.}) Insekto di qua la ali esas membrana (abeli,}
\l{\cel{3}vespi, formiki, e c.) - dEFIRS.}
\l{}
\l{\sou{himno}. I. Poemo religiala, honore di la dei, di la heroi. - II.}
\l{\cel{3}Poemo lirika, honore di la deajo. - III. Kantiko Latina quan}
\l{\cel{3}onu kantas o recitas en kirko. - DEFIRS.}
\l{}
\l{\sou{hioido}. (\sou{anat.}) (Osto) qua esas horizontale inter la bazo di la}
\l{\cel{3}lango e la laringo. - EFIS.}
\l{}
\l{\sou{hipar}. (\sou{netrans}.)Sentar ta kontrakto intervalopa qua sukusas}
\l{\cel{3}la diafragmo ed akompanesas da bruiso ne-artikulata partiku-}
\l{\cel{3}lara. - S.}
\l{}
\l{\sou{hipalago}. (\sou{retor.}) Stilo-figuro qua atribuas a vorto to quo}
\l{\cel{3}konvenas ad altra vorto. (kom ex.: retrodonar ulu a vivo,}
\l{\cel{3}vice : retrodonar vivo ad ulu). - DEFIS.}
\l{}
\l{hiper-. Prefixo qua signifikas \textquotedbl{}super, trans, ecese\textquotedbl{}.}
\l{}
\l{hiperbato. (\sou{retor.}) Figuro qua konsistas ye permutar la ordino}
\l{\cel{3}naturala di ula vorti en propoziciono. - DEFIS.}
}
